\documentclass{situationist-preprint}

\title{Divine Command and Joint Grounding}
\msid{MTE-20-a}
\msfamily{MTE-20}
\msdiscipline{Metaethics}
\mstype{Principal manuscript}
\mskeywords{divine command theory; grounding; moral obligation; co-endorsement; Euthyphro dilemma; theistic metaethics}

\begin{document}
\maketitle

\begin{abstract}
Does God ground morality? Divine command theorists hold that God's commands make acts obligatory; theological voluntarists hold more generally that divine facts are the full ground of moral facts. We argue that, on theism, both theses are false, and that a third is true: morality is jointly grounded in divine and non-divine facts. The argument applies the co-endorsement account of REL-09-a to moral obligation. An act-type is obligatory in virtue of being co-endorsed, as required, by God and by the natures of the creatures it concerns; co-endorsement is grounded in a divine fact and in a non-divine fact, and the non-divine fact is not itself grounded in divine facts. Divine commands are not among the grounds of obligation: God commands only what is co-endorsed, so that a command guarantees and makes known the obligation it concerns without creating it. Divine commands confirm rather than create moral obligations. The answer to the question is therefore no, in the sense in which divine command theory says yes.
\end{abstract}

\section{The question}

The claim that God grounds morality has two classical forms. In the voluntarist tradition associated with Ockham, moral facts are constituted by the divine will, and what God wills could have been otherwise. In contemporary divine command theory the claim is restricted to obligation: \citet{Quinn1978} held that moral requirements depend on divine commands, and \citet{Adams1999} and \citet{Evans2013} that obligation is constituted by the commands of a good God, while goodness is grounded in God's nature. The objections are equally classical. If commands make obligations, then had God commanded cruelty, cruelty would have been obligatory; if they do not, God's commands seem idle.

REL-09-a resolves the Euthyphro dilemma by denying its presupposition that one of goodness and divine command grounds the other. Both are grounded in co-endorsement: an act-type is co-endorsed when it is endorsed both by God and by the natures of the creatures it concerns (REL-09-a, Definition~2.1). Treating grounding as irreflexive and transitive (Assumption~2.4) and co-endorsement as the full common ground of goodness and of command (Assumptions~2.6 and~2.7), REL-09-a proves that neither horn holds (Theorem~3.2), that goodness is grounded in both a divine and a non-divine fact (Proposition~3.3), and that if God commands an act-type it is good, although the command does not ground its goodness (Corollary~3.4). REL-09-a notes that the argument applies without change to the fact that an act-type is required. The present manuscript draws the metaethical consequence for obligation and answers the question in the title.

\section{From goodness to obligation}

Read \enquote{required} for \enquote{good} throughout the framework of REL-09-a. An act-type $a$ is \emph{co-endorsed as required} when God, from God's own nature, wills that $a$ be done, and the natures of the creatures $a$ concerns are such that their flourishing calls for $a$ to be done. The assumptions become:

\begin{assumption}\label{a:transfer}
The fact that $a$ is co-endorsed as required is grounded in the fact that God endorses $a$ as required and in the fact that creaturely natures call for $a$, and obtains just in case both do. It grounds both the fact that $a$ is obligatory and the fact that God commands $a$; every ground of either fact is that co-endorsement or a ground of it; and $a$ is obligatory just in case it is co-endorsed as required, while God commands $a$ only if it is.
\end{assumption}

\begin{assumption}\label{a:independent}
The fact that creaturely natures call for $a$ is not grounded in any divine fact.
\end{assumption}

Assumption~\ref{a:transfer} is the conjunction of REL-09-a's Assumptions~2.5 to~2.7 and~2.9 with \enquote{obligatory} in place of \enquote{good}; Assumption~2.4, on grounding, and Assumption~2.8, on which of the two endorsements is divine, carry over unchanged. Assumption~\ref{a:independent} is the claim REL-09-a defends in reply to the objection that the dilemma recurs: God chooses which natures to instantiate, but not what a creature of a given nature needs in order to flourish, which is fixed by the essence of its kind, as MET-13 argues that essence is prior to modality.

\begin{proposition}\label{p:joint}
For every act-type $a$, the fact that $a$ is obligatory is grounded in a divine fact and in a non-divine fact, and $a$ is obligatory only if both obtain.
\end{proposition}

\begin{proof}
REL-09-a, Proposition~3.3, under Assumption~\ref{a:transfer}.
\end{proof}

\begin{proposition}\label{p:notfull}
Divine facts are not the full ground of any obligation.
\end{proposition}

\begin{proof}
By Proposition~\ref{p:joint}, the obligation of $a$ obtains only if creaturely natures call for $a$, and the fact that they do is among its grounds. By Assumption~\ref{a:independent}, that fact is not grounded in any divine fact. So the divine facts do not by themselves ground the obligation: a ground of it lies outside them and is not grounded in them.
\end{proof}

\begin{proposition}\label{p:confirm}
If God commands $a$, then $a$ is obligatory; but God's command of $a$ does not ground the obligation of $a$.
\end{proposition}

\begin{proof}
REL-09-a, Corollary~3.4, under Assumption~\ref{a:transfer}. The first clause follows because God commands only what is co-endorsed as required, and what is so co-endorsed is obligatory; the second because the command and the obligation have co-endorsement as their common ground, and by REL-09-a, Theorem~3.2, neither grounds the other.
\end{proof}

\begin{theorem}\label{t:main}
God does not ground morality, either in the sense that divine facts are its full ground or in the sense that divine commands create obligations. Morality is jointly grounded in divine and non-divine facts, and divine commands confirm rather than create moral obligations.
\end{theorem}

\begin{proof}
Propositions~\ref{p:joint}, \ref{p:notfull} and~\ref{p:confirm}.
\end{proof}

A command \emph{confirms} an obligation in two senses. It guarantees it, since God commands only what is obligatory; and it makes it known and expresses God's endorsement, which is one of the obligation's grounds. What it does not do is add a ground. The divine fact that does ground obligation is God's endorsement, not God's command, and the distinction is the point at which the present view departs from divine command theory.

\section{Objections}

\begin{objection}[Obligation needs a commander]
Adams argues that obligation, unlike goodness, is a social notion: to be obligated is to be subject to a legitimate demand, and demands must be made \citep{Adams1999}. Without a command there is goodness but no obligation.
\end{objection}

\begin{reply}
The social element Adams identifies is present in co-endorsement. God's endorsement of $a$ as required is the will of a person, who loves the creatures $a$ concerns, that $a$ be done; it is a demand in every respect but that of being communicated. The command communicates it. What the objection shows is that obligation involves a person's will, which the present view accepts; it does not show that it involves the act of commanding over and above the will the command expresses.
\end{reply}

\begin{objection}[The answer should be yes]
On the view, nothing is obligatory that God does not endorse, and God's will is among the grounds of every obligation. That is enough for God to ground morality.
\end{objection}

\begin{reply}
If grounding morality means being among its grounds, the view agrees, and Proposition~\ref{p:joint} says so. The question as it has been debated, between divine command theorists and their critics, is whether God's will or command is what makes acts right, the first horn of the Euthyphro dilemma. In that sense the answer is no: divine facts are a partial ground and commands are no ground at all.
\end{reply}

\begin{objection}[Idle commands]
If commands create no obligations, they serve no purpose.
\end{objection}

\begin{reply}
They serve the purpose of every reliable testimony to a fact one cannot easily discover for oneself, and they express the endorsement of the one whose will is among the grounds of the fact. A command is redundant as a ground, not as a communication, as REL-09-a also replies.
\end{reply}

\section{Conclusion}

Applied to obligation, the co-endorsement account of REL-09-a implies that obligations are grounded jointly in God's endorsement and in what creaturely natures call for, that the second ground is not itself divine, and that God's commands follow and guarantee obligations without grounding them. God does not ground morality in the sense in which divine command theory holds that God does: morality is jointly grounded in divine and non-divine facts, and divine commands confirm rather than create moral obligations.

\begin{thebibliography}{99}\small
\bibitem[Adams(1999)]{Adams1999} Adams, R.~M. (1999). \emph{Finite and Infinite Goods: A Framework for Ethics}. Oxford University Press.
\bibitem[Evans(2013)]{Evans2013} Evans, C.~S. (2013). \emph{God and Moral Obligation}. Oxford University Press.
\bibitem[Quinn(1978)]{Quinn1978} Quinn, P.~L. (1978). \emph{Divine Commands and Moral Requirements}. Clarendon Press.
\end{thebibliography}

\end{document}
